\subsection{Die schriftliche Ausarbeitung}



\subsubsection{Umfang}
Umfang: so kurz wie möglich und so umfangreich wie nötig (in Absprache mit dem/der Betreuer/in)

Die schriftliche Abschlussarbeit soll eine geordnete Darstellung der im Verlauf der Arbeit angestellten Überlegungen, durchgeführten Ableitungen, entwickelten Algorithmen, Experimenten usw. im Sinne einer Diskussion bringen. Das Herausarbeiten und Bewerten der Ergebnisse ist dabei die wesentliche geistige Leistung. Eine konzentrierte Darstellung (etwa $30$ bis $50$ A4- Seiten) ist anzustreben; vermeide unbedingt weitschweifige Abhandlungen. Experimentprotokolle, Rechnerausdrucke, etc. können geordnet, eindeutig und vollständig beschriftet als Anhang beigeheftet wer-den. Berechnungen sind trotz der gebotenen Kürze so ausführlich zu bringen, sodass eine schrittweise Überprüfung möglich ist.

\subsubsection{Äußere Form}

Die Arbeit ist in Deutsch oder Englisch abzufassen und mit dem Textverarbeitungssystem LaTeX oder doc zu schreiben. Die bereitgestellte Vorlage muss genutz werden.
Eine Tex-Vorlage in OVERLEAF!
\hyperlink{https://www.overleaf.com/project/629f45ce28c12f2c468b8f78}{
Vorlage für Abschlussarbeit an der HTW}

\begin{itemize}
\item  Verwende niemals fremde Ideen/ Texte/ Bilder, ohne die Urheberrechte zu nennen und zu klären. 
\item  Analyze the papers you read with respect to their style.	
\item  Beginne mit einer logischen Struktur, exakte Formulierungen mache später. Mehr dazu in Kap.~\ref{sec:struktur}
\item Lass dir deine Arbeit von jemandem überarbeiten.
\item  Die Zusammenfassung sollte besonders gut werden. Viele Leser werden nur diese lesen.
\item  Auch wenn für die endgültige Fassung kein Inhaltsverzeichnis benötigt wird, macht es Sinn, mit Hilfe des Befehls 
tableofcontents während des Schreibens eines zu erstellen. Damit behalten Sie Ihre Struktur im Auge. % Use the command \tableofcontents}for that purpose. 
\item Bemühen Sie sich um Symmetrie in Bezug auf die Anzahl und Länge der Abschnitte pro Kapitel oder der Unterabschnitte pro Abschnitt (Einleitung/ Schluss sind Ausnahmen mit weniger Unterüberschriften). 



\item Zwischen Abschnittsüberschriften und Unterabschnittsüberschriften sollten keine "losen" Textabschnitte eingefügt werden. Ein oder zwei Sätze sind möglich, aber nur für einen Überblick über den gesamten Abschnitt.


\item Ein Abschnitt kann nicht nur einen einzigen Unterabschnitt enthalten.

\item Überspringen Sie keine Überschriftenebenen.


\item Vermeiden Sie mehr als drei Ebenen (Abschnitt, Unterabschnitt, Unterunterabschnitt). Wenn Sie auf einer tieferen Ebene unterteilen müssen, überarbeiten Sie Ihre Struktur.

\item Jeder Absatz sollte eine eindeutige Idee/ Botschaft enthalten. Außerdem sollte der Leser das Thema des Absatzes bereits in den ersten drei Wörtern, mindestens aber im ersten Satz erkennen.

\item Achten Sie auf eine logische Verbindung innerhalb und zwischen den Absätzen. Am einfachsten ist es, zuerst die Hauptgedanken aufzuschreiben und sie erst dann zu Absätzen auszubauen.

\item  Do not get creative with different types of line breaks between paragraphs. If you find that necessary, revise your structure.

\item Die Leser sollten schnell wissen, wo sie welche Informationen finden können. Deshalb ist wichtig, klare Überschriften zu verwenden. Keine Abkürzungen in Überschriften.

\end{itemize}